\rhead{CY-EL4: Blockchain \& Decentralized Identity}
\phantomsection 
\section*{Blockchain \& Decentralized Identity (CY-EL4)}
\label{elec:CY_4}

\noindent \small{\hyperref[main:scheme]{$\leftarrow$ Back to Scheme}} \vspace{0.5cm}

\noindent \textbf{Credits:} 4 (3-0-2)

\subsection*{Theory}
\noindent\textbf{Unit 1} \\
\textit{Blockchain Fundamentals and Cryptography:} Distributed ledger technology principles, Cryptographic hash functions (SHA-256, Keccak), Digital signatures (ECDSA, EdDSA), Merkle trees and Patricia tries, Consensus mechanisms (PoW, PoS, PBFT, Tendermint), Byzantine Generals Problem, 51\% attacks and chain reorganizations, Blockchain trilemma (decentralization, security, scalability).

\vspace{0.3cm}

\noindent\textbf{Unit 2} \\
\textit{Smart Contracts and Ethereum Ecosystem:} EVM architecture and gas model, Solidity programming (contracts, modifiers, events), Smart contract vulnerabilities (reentrancy, integer overflow, front-running), Formal verification (Mythril, Slither, Certora), ERC standards (20, 721, 1155, 4337 account abstraction), Layer 2 scaling (Optimism, Arbitrum, zkSync), Oracles (Chainlink, Band Protocol).

\vspace{0.3cm}

\noindent\textbf{Unit 3} \\
\textit{Decentralized Identity (DID) Standards:} Self-Sovereign Identity (SSI) principles, DID method specification (did:ethr, did:key, did:web), Verifiable Credentials (VC) and Presentations (VP) data model, W3C standards (DID Core, VC Data Model), Decentralized Identifiers resolution, Issuer/Holder/Verifier model, Zero-knowledge proofs for selective disclosure.

\vspace{0.3cm}

\noindent\textbf{Unit 4} \\
\textit{DID Ecosystems and Interoperability:} DIDComm protocol for secure messaging, DIF standards (Universal Resolver, Sidetree protocol), Hyperledger Indy/Aries for permissioned identity, uPort/SeraphId implementations, Verifiable Data Registries (VDRs), Revocation and expiration handling (status registries), Cross-chain identity (Polkadot, Cosmos IBC).

\vspace{0.3cm}

\noindent\textbf{Unit 5} \\
\textit{Blockchain Security and Privacy:} 51\% attacks mitigation, Eclipse attacks, Smart contract auditing methodologies, Flash loan attacks and MEV, Privacy-preserving techniques (zk-SNARKs, Bulletproofs, Mimblewimble), Confidential transactions, Threshold signatures and MPC wallets, Quantum threats (lattice-based crypto, hash-based signatures), Regulatory compliance (MiCA, travel rule).

\newpage

\subsection*{Practical}
\noindent\textbf{Lab 1: Blockchain Cryptography Basics} \\
Focuses on ECDSA signing/verification, Merkle tree construction, transaction serialization.

\vspace{0.2cm}

\noindent\textbf{Lab 2: Ethereum Smart Contracts} \\
Focuses on Remix IDE development, Hardhat testing, gas optimization.

\vspace{0.2cm}

\noindent\textbf{Lab 3: Smart Contract Vulnerability Exploitation} \\
Focuses on reentrancy demos, Slither static analysis, remediation.

\vspace{0.2cm}

\noindent\textbf{Lab 4: DID Creation and Resolution} \\
Focuses on did:key/ethr generation, Universal Resolver usage.

\vspace{0.2cm}

\noindent\textbf{Lab 5: Verifiable Credentials Issuance} \\
Focuses on VC signing, VP presentation, selective disclosure.

\vspace{0.2cm}

\noindent\textbf{Lab 6: DIDComm Secure Messaging} \\
Focuses on Aries agents, pairwise DIDs, encrypted messages.

\vspace{0.2cm}

\noindent\textbf{Lab 7: Layer 2 Deployment} \\
Focuses on Optimism/Arbitrum bridges, L2 contract deployment.

\vspace{0.2cm}

\noindent\textbf{Lab 8: zk-SNARK Implementation} \\
Focuses on Semaphore/circom circuits, proof generation/verification.

\vspace{0.2cm}

\noindent\textbf{Lab 9: Blockchain Forensics} \\
Focuses on transaction tracing, tainted address analysis.

\vspace{0.2cm}

\noindent\textbf{Lab 10: SSI Application Prototype} \\
Focuses on complete KYC/credential verification dApp.

\vspace{0.2cm}

\newpage
\rhead{Curriculum Schema}
